\documentclass[a5paper,10pt]{article}

% https://github.com/InDevRus/filippov-solutions
% Нумерация страниц
\pagenumbering{gobble}

% Настройка полей
\usepackage{geometry}
\geometry{left=1cm}
\geometry{right=1cm}
\geometry{top=1cm}
\geometry{bottom=1.5cm}

% Вставка таблицы
\usepackage{array}
\usepackage{tabularx}

% Инструменты выравнивания формул
\usepackage{amsmath}
\usepackage{mathtools}

% Диакритика в математике
\usepackage{amsfonts}

% Вычисления размеров на ходу
\usepackage{calc}

% Удобные записи производных
\usepackage[italic=true]{derivative}

% Настройки сносок
\usepackage{enotez}

% Длинные знаки векторов
\usepackage{anyfontsize}
\usepackage[b]{esvect}

% Вставка картинок
\usepackage{graphicx}

% Вставка красной строки
\usepackage{indentfirst}
\setlength{\parindent}{1em}

% Условия в командах
\usepackage{xifthen}

% Луа-скрипты
\usepackage{luacode}

% Настройка команд отдельно в математическом и текстовом режимах
\usepackage{mathcommand}

% Для повторения LaTeX кода
\usepackage{multido}

% Конфигурация отступов формул
\delimitershortfall-1sp
\usepackage{mleftright}
\mleftright

% Растягивание объектов
\usepackage{relsize}

% Обтекание текстом
\usepackage{wrapfig}
\usepackage{subcaption}
\usepackage{floatflt}

% Поддержка ввода кириллицы
\usepackage[english, russian]{babel}

% Настройка корневой позиции для компилятора
\usepackage{import}
\providecommand*{\ProjectRootPath}{..}

\import{\ProjectRootPath/preambles/macros/}{theorems}

\import{\ProjectRootPath/preambles/fonts}{font_setup}

% Знак градуса
\usepackage{siunitx}

\import{\ProjectRootPath/preambles/macros/}{operators}
\import{\ProjectRootPath/preambles/macros/}{redefinitions}

\import{\ProjectRootPath/preambles/fonts}{letters_setup}
\import{\ProjectRootPath/preambles/fonts/adjustments}{custom_superscripts}

\import{\ProjectRootPath/preambles/custom}{formatting}
\import{\ProjectRootPath/preambles/custom}{frames}
\import{\ProjectRootPath/preambles/custom}{list_setup}

% TODO: перевести команды на современный стиль объявления

\begin{document}

\begin{framed}
    Если в полярных координатах $\tsys{& x = r \cos\theta \\& y = r \sin\theta}$ имеем \vspace{1mm} производную $\dfrac {dr} {d\theta} = A\br{r; \theta}$,
    то соответствующая производная в декартовых координатах $\dfrac {dy} {dx} = B\br{r; \theta}$ связана с исходной $A\br{r; \theta}$ соотношениями
    \begin{align*}
        \dfrac {dy} {dx}
        = B\br{A; r; \theta}
        = \dfrac {A \sin\theta + r \cos\theta} {A \cos\theta - r \sin\theta};&&
        \dfrac {dr} {d\theta}
        = A\br{B; r; \theta}
        = r \cdot \dfrac {B \sin\theta + \cos\theta} {B\cos\theta - \sin\theta}.
    \end{align*}
\end{framed}

\begin{framed}
    Для угла $\phi = \ang{90}$ имеем $\tg{\beta} = \tg\br{\alpha \pm \phi} =\tg\br{\alpha \pm \ang{90}} = -\ctg \alpha = -\dfrac {1} {\tg\alpha}$.

    С учётом указанного выше для данного угла $\phi$ всякое семейство изогональных траекторий $z\br{x}$ и соответствующее семейство в полярной системе $\tsys{& x = \rho \cos\theta \\& z = \rho \sin\theta}$ удовлетворяют соотношениям
    \begin{align*}
        A_{z}
        = \rho \cdot \dfrac {B_{z} \sin\theta + \cos\theta} {B_{z}\cos\theta - \sin\theta};
        &&
        B_{z}\br{\rho; \theta} = -\dfrac {1} {B_{y}\br{\rho; \theta}};
        &&
        B_{y} = \dfrac {A_{y} \sin\theta + r \cos\theta} {A_{y} \cos\theta - r \sin\theta},
    \end{align*}
    где введены обозначения
    $A_{y} = \dfrac {dr} {d\theta}$,
    $B_{y} = \dfrac {dy} {dx}$,
    $A_{z} = \dfrac {d\rho} {d\theta}$,
    $B_{z} = \dfrac {dz} {dx}$. \vspace{1mm}

    Последовательно подставив и затем упростив, получим следующую формулу:
    $$A_{z}\br{\rho; \theta} = -\dfrac {\rh\pow{o}{2}} {A_{y}\br{\rho; \theta}}.$$
\end{framed}

$r = a + \cos\theta$.
$A_{y} = r\'\br{\theta} = -\sin\theta$.
Таким образом, уравнение изогональных (ортогональных) траекторий в обозначениях, указанных выше, принимает форму
$\dfrac {d\rho} {d\theta} = A_{z} = \linebreak = -\dfrac {\rh\pow{o}{2}} {A_{y}} = \dfrac {\rh\pow{o}{2}} {\sin\theta}.$
$\sin\theta\ \rho\'\br{\theta} = \rh\pow{o}{2}\br{\theta}$.

\begin{remark}
    Уравнение
    $\rho\' = \dfrac {\rh\pow{o}{2}} {\sin\theta}$
    имеет общее решение
    $$\rho = \dfrac {1} {C - \ln\vbr{\tg\dfrac {\theta} {2}}}.$$
\end{remark}

\end{document}

